% ============================================================
% Section 7: Conclusion
% The Information Budget of Graph Neural Networks
% ============================================================

\section{Conclusion}
\label{sec:conclusion}

We present the \textbf{Information Budget} framework, a unified theoretical perspective that explains when GNNs improve upon feature-only methods and when they cause harm. The key insight is that MLP accuracy determines the \textit{ceiling} for GNN improvement, while homophily determines the \textit{utilization} of this ceiling.

\subsection{Summary of Contributions}

\textbf{1. The Conditional Enhancement Bound}: We establish that GNN improvement is upper-bounded by $\mathcal{B} = 1 - \text{Acc}_{\text{MLP}}$ under fair comparison conditions (equivalent model capacity, identical features, standard aggregation). This resolves the Same-$h$ Puzzle: Cora ($h=0.81$, Budget=24.9\%) gains +12.7\% from GCN, while Coauthor-CS ($h=0.81$, Budget=5.6\%) gains $-0.7\%$. Same homophily, opposite outcomes, different budgets. Validated on:
\begin{itemize}
    \item 9 noise levels on Cora: All within budget bound (9/9)
    \item 19 benchmark datasets: All satisfy the constraint (19/19)
    \item Edge shuffle experiment: Causal validation (47.4\% drop when structure randomized)
    \item 4 OGB datasets at million-node scale: All satisfy the constraint
\end{itemize}

\textbf{2. Aggregation Damage Ratio (ADR)}: We quantify how aggregation can \textit{destroy} information:
\begin{equation}
    \text{ADR} = 1 - \frac{\text{Acc}_{\text{GNN}}}{\text{Acc}_{\text{MLP}}}
\end{equation}
Key findings:
\begin{itemize}
    \item ADR peaks at $h = 0.5$ (mid-homophily): GCN loses 18.9\% of MLP's information
    \item Replace architectures (GCN, GAT): ADR = 0.05--0.19 at mid-$h$
    \item Concatenate architectures (GraphSAGE): ADR $\approx$ 0.001 in controlled settings ($\sim$189$\times$ reduction), $\sim$20$\times$ reduction on real data
    \item \textbf{Information-theoretic explanation}: By the Data Processing Inequality, concatenation preserves $I(Y; X, \tilde{X}) \geq I(Y; X)$, while replacement suffers from Aggregation Information Loss
    \item Ablation confirms: Adding concatenation to GCN recovers performance ($-24.6\% \to +3.1\%$)
\end{itemize}

\textbf{3. Dual Heterophily Types}: We identify two fundamentally different low-$h$ regimes:
\begin{itemize}
    \item \textbf{Type A} (Recoverable): $R_{2\text{-hop}} > 1.5\times$. H2GCN improves +28.1\% average. Examples: Texas, Wisconsin, Cornell.
    \item \textbf{Type B} (Irrecoverable): $R_{2\text{-hop}} \leq 1.0\times$. MLP is optimal unless features are weak. Examples: Actor, Chameleon, Squirrel.
\end{itemize}
This resolves the Chameleon-Squirrel paradox: they are Type B with unusually weak features (MLP = 32--49\%), so ``bad information $>$ no information'' applies.

\textbf{4. Unified Decision Framework}: Algorithm~\ref{alg:budget_decision} combines Budget, ADR, and Heterophily Type:
\begin{itemize}
    \item CSBM controlled experiments: 88.9\% prediction accuracy (32/36)
    \item External validation: 77.8\% accuracy (7/9 datasets)
    \item Key insight: Check budget first, then $h$, then heterophily type
\end{itemize}

\subsection{The Paradigm Shift}

This work represents a shift in understanding GNN behavior:

\begin{center}
\begin{tabular}{@{}l|l@{}}
\textbf{Old View} & \textbf{New View} \\
\midrule
``High $h$ $\Rightarrow$ GNN wins'' & ``High $h$ + sufficient budget $\Rightarrow$ GNN may win'' \\
``Low $h$ $\Rightarrow$ GNN loses'' & ``Low $h$ Type A $\Rightarrow$ H2GCN; Type B $\Rightarrow$ depends on budget'' \\
``GNN $>$ MLP always'' & ``GNN can damage information; MLP is a strong baseline'' \\
``More layers = better'' & ``More aggregation = more potential damage'' \\
\end{tabular}
\end{center}

\subsection{Practical Guidelines}

\textbf{For Practitioners}:
\begin{enumerate}
    \item \textbf{Always train MLP first}: This reveals the budget. If $\text{Acc}_{\text{MLP}} > 0.95$, stop---no room for GNN improvement.

    \item \textbf{Use concatenation when uncertain}: GraphSAGE or GCN-Concat provides robustness against aggregation damage.

    \item \textbf{For low-$h$ datasets}: Compute $R_{2\text{-hop}}$. If $> 1.5$, use H2GCN/LINKX. If $\leq 1.0$ with high MLP, use MLP.

    \item \textbf{On U-shape patterns}: The symmetric U-shape (GNN performance dipping at $h \approx 0.5$) is observed in controlled synthetic experiments (CSBM) where feature quality is held constant. In real datasets, confounding factors (varying feature quality, degree distributions) often obscure this pattern. Focus on computing $\kappa = \rho^2 k$ rather than expecting a visual U-shape.
\end{enumerate}

\textbf{For Researchers}:
\begin{enumerate}
    \item Report MLP baseline and budget in all GNN papers
    \item Compare against MLP, not just other GNNs
    \item Analyze failure cases through the ADR lens
    \item Distinguish Type A vs Type B when studying heterophily
\end{enumerate}

\subsection{Broader Impact}

\textbf{Scientific Contribution}: The Information Budget framework shifts focus from ``proposing better architectures'' to ``understanding when architectures work.'' This diagnostic approach can prevent wasteful experimentation and guide principled method selection.

\textbf{Practical Value}: By revealing the budget constraint upfront, practitioners can avoid overengineering on problems where features are already sufficient. This saves computational resources and reduces deployment complexity.

\textbf{Generalization}: While validated on node classification, the principles generalize:
\begin{itemize}
    \item Link prediction: Budget = $1 - \text{Acc}_{\text{feature-only}}$
    \item Graph classification: Budget = $1 - \text{Acc}_{\text{global pooling}}$
    \item Heterogeneous graphs: Budget varies by node/edge type
\end{itemize}

\subsection{Limitations and Future Work}

\textbf{1. Feature Quality Measurement}: We use MLP accuracy as a proxy for feature sufficiency. Direct feature quality metrics could provide faster diagnostics.

\textbf{2. Dynamic Graphs}: The budget may change as graphs evolve. Temporal analysis of budget dynamics is future work.

\textbf{3. Multi-Relational Graphs}: Different edge types may have different budgets and ADRs. Extending the framework to knowledge graphs is promising.

\textbf{4. Learned Heterophily Type}: Automatic classification of Type A vs Type B without computing 2-hop homophily could accelerate method selection.

\subsection{Closing Remarks}

The GNN community has produced hundreds of architectures, each claiming improvements on specific benchmarks. What has been missing is a principled understanding of \textit{when} and \textit{why} each approach works. The Information Budget framework fills this gap.

Our key message: \textbf{Features determine the ceiling; structure determines the path to it; aggregation can block the path entirely.}

When encountering a new graph learning task, the first question should not be ``which GNN to use?'' but rather ``\textbf{is there budget for a GNN?}'' This simple diagnostic---training an MLP first---can save significant effort and lead to more appropriate architectural choices.

\textbf{Reproducibility}: Code, experiments, and metric computation tools are available at \url{https://github.com/MengdanXue/trust-regions-gnn}. All experiments use 10 random seeds with statistical significance testing.
